\documentclass{article}

\usepackage{amsmath,amssymb}
\usepackage{xurl}
\usepackage[hidelinks]{hyperref}
\setlength{\emergencystretch}{2em}

% Shared commands for notes in docs/paper-gaps/.

\DeclareUnicodeCharacter{2080}{\ifmmode{}_{0}\else\textsubscript{0}\fi}
\DeclareUnicodeCharacter{2081}{\ifmmode{}_{1}\else\textsubscript{1}\fi}
\DeclareUnicodeCharacter{2082}{\ifmmode{}_{2}\else\textsubscript{2}\fi}
\DeclareUnicodeCharacter{2099}{\ifmmode{}_{n}\else\textsubscript{n}\fi}

\newcommand{\Lean}{\textsc{Lean}}
\ExplSyntaxOn
\NewDocumentCommand{\leanid}{m}{
  \ifmmode
    \textsf{\detokenize{#1}}
  \else
    \group_begin:
    \urlstyle{sf}
    \tl_set:Nn \l_tmpa_tl {#1}
    \tl_replace_all:Nnn \l_tmpa_tl {\_} {_}
    \exp_args:NV \path \l_tmpa_tl
    \group_end:
  \fi
}
\ExplSyntaxOff
\providecommand{\tr}{\operatorname{tr}}
\newcommand{\tnleanIssueBase}{https://github.com/LionSR/TNLean/issues/}
\newcommand{\tnleanPrBase}{https://github.com/LionSR/TNLean/pull/}
\newcommand{\tnleanDocBase}{https://sirui-lu.com/TNLean/docs/}
\newcommand{\ghissue}[1]{\href{\tnleanIssueBase#1}{GitHub issue \##1}}
\newcommand{\ghpr}[1]{\href{\tnleanPrBase#1}{PR \##1}}
\newcommand{\doclink}[1]{\href{\tnleanDocBase#1.html}{\path{#1}}}

% Dirac notation and the identity operator, as in blueprint/src/macros/common.tex.
% \providecommand keeps note-local definitions authoritative.
\usepackage{dsfont}
\providecommand{\ket}[1]{|#1\rangle}
\providecommand{\bra}[1]{\langle #1|}
\providecommand{\braket}[2]{\langle #1 | #2 \rangle}
\providecommand{\ketbra}[2]{|#1\rangle\!\langle #2|}
\providecommand{\Id}{\mathds{1}}
\providecommand{\one}{\mathds{1}}
\providecommand{\C}{\mathbb{C}}
\providecommand{\MN}[1]{M_{#1}(\C)}
\providecommand{\MD}{M_D(\C)}

% External referencing for paper sources.  The build helper writes sanitized
% auxiliary files under build/paper-aux/ and excludes duplicated source labels.
\usepackage{xr}

% Import a generated paper auxiliary file with a caller-supplied label prefix.
% The candidate paths support builds from docs/paper-gaps/, the repository root,
% and build/paper-aux/ itself.
\newcommand{\importpaperaux}[2]{%
  \IfFileExists{../../build/paper-aux/#2.aux}{%
    \externaldocument[#1]{../../build/paper-aux/#2}%
  }{%
    \IfFileExists{build/paper-aux/#2.aux}{%
      \externaldocument[#1]{build/paper-aux/#2}%
    }{%
      \IfFileExists{#2.aux}{%
        \externaldocument[#1]{#2}%
      }{}%
    }%
  }%
}

% General external reference macros
\newcommand{\paperref}[2]{\ref{#1-#2}}
\newcommand{\papereqref}[2]{(\ref{#1-#2})}

% CPSV16 source (1606.00608 / MPDO-22-12-17-2.tex) external document and shortcuts
\importpaperaux{cpsv16-}{1606.00608/MPDO-22-12-17-2-tnlean}
\newcommand{\cpsvref}[1]{\paperref{cpsv16}{#1}}
\newcommand{\cpsveqref}[1]{\papereqref{cpsv16}{#1}}

% Verdict marker: \gapnote{<kind>}{<status>}. Kind is the policy
% classification (clarification, local-correction, scope-restriction,
% unfaithful, false-source, open-gap); status is open, wip (elimination
% underway), resolved, or historical. Machine-read by the site index;
% prints a verdict line.
\newcommand{\gapnote}[2]{\par\noindent\textbf{Verdict:} #1 (#2).\par}

\providecommand{\leanid}[1]{\path{#1}}

\title{Injective PEPS Fundamental Theorem: Section 3 Route}
\author{}
\date{2026-05-06}

\begin{document}
\maketitle
\gapnote{clarification}{resolved}

\section{The source argument}

Section~3 of Ref.~\cite{Molnar2018NormalPEPS} proves the injective projected
entangled pair state (PEPS) Fundamental Theorem for finite graphs without
double edges or self-loops. The proof is local: it fixes an edge, blocks the
graph into a three-site injective matrix product state (MPS), applies the
one-dimensional injective Fundamental Theorem, and repeats the construction at
every edge. It does not reduce a two-dimensional network to a long
one-dimensional chain.\footnote{The edge-blocking construction is
\path{Papers/1804.04964/paper_normal.tex:981--1009}. The equality of the two
PEPS before blocking is at
\path{Papers/1804.04964/paper_normal.tex:1010--1036}. The application of the
three-site injective MPS theorem is at
\path{Papers/1804.04964/paper_normal.tex:1037--1065}.}

Fix an edge $e=(u,v)$. The endpoint tensors at $u$ and $v$ form the first and
third MPS sites, and the contraction of all remaining vertices forms the middle
site. Injectivity of the PEPS tensors implies injectivity of this blocked MPS.
Because the two original PEPS represent the same state, their blocked MPS also
represent the same state. The one-dimensional injective Fundamental Theorem
then gives an algebra isomorphism between matrix insertions on the two copies of
the chosen bond. Thus, for every matrix $X$ inserted on that bond in the first
network, a matrix on the corresponding bond in the second network produces the
same state.

Choose these edge gauges for every edge and absorb them into the second tensor
family. Equation \path{eq:inj_equal_edge} of
Ref.~\cite{Molnar2018NormalPEPS} then gives a stronger comparison: for every
edge and every matrix $X$, inserting $X$ on that edge produces the same state
from the first PEPS and the modified second PEPS.\footnote{The relevant display
is Equation \path{eq:inj_equal_edge} in
\path{Papers/1804.04964/paper_normal.tex:1037--1065}.}

\section{Formal representation}

The definitions in \path{TNLean/PEPS/Defs.lean} represent finite graphs, edges,
incident edges, tensors, virtual configurations, state coefficients, equality
of states, and vertex injectivity. Vertex injectivity means that the local
physical vectors indexed by local virtual configurations are linearly
independent. This agrees with the injectivity assumption in
Ref.~\cite{Molnar2018NormalPEPS}, where each tensor is injective as a map from
its virtual space to its physical space.

The file \path{TNLean/PEPS/VirtualInsertion.lean} realizes virtual operations
at an injective vertex. A chosen left inverse of the local tensor map converts
a local virtual endomorphism into a physical operation. This is the local
linear-algebraic step used in the three-site MPS argument of
Ref.~\cite{Molnar2018NormalPEPS}. The basis identity
\leanid{TNLean.PEPS.localIncidentMatrixOp_single} and its tensor-map form
\leanid{TNLean.PEPS.localTensorMap_localIncidentMatrixOp_single} identify the
matrix action on one virtual boundary leg with the sum over the new index on
the distinguished edge.

The file \path{TNLean/PEPS/Blocking.lean} contains the edge-centered blocking
data. Its declarations split a local virtual configuration at a distinguished
incident edge and decompose the vertex product into the left endpoint, the
middle region, and the right endpoint. The edge-blocked coefficient is proved
equal to the original PEPS coefficient.

The matrix-insertion layer uses
\leanid{TNLean.PEPS.EdgeInsertedBoundaryConfig},
\leanid{TNLean.PEPS.edgeBoundaryToInsertedBoundaryConfig},
\leanid{TNLean.PEPS.EdgeDiagonalInsertedBoundaryConfig},
\leanid{TNLean.PEPS.edgeBoundaryConfigEquivDiagonalInsertedBoundaryConfig},
\leanid{TNLean.PEPS.edgeMiddleConfigToOpenMiddleConfig},
\leanid{TNLean.PEPS.edgeOpenMiddleConfigToMiddleConfig},
\leanid{TNLean.PEPS.edgeMiddleConfigEquivOpenMiddleConfig},
\leanid{TNLean.PEPS.edgeOpenMiddleWeight},
\leanid{TNLean.PEPS.edgeBoundaryLeftIndexEquivInsertedBoundaryConfig},
\leanid{TNLean.PEPS.edgeBoundaryRightIndexEquivInsertedBoundaryConfig}, and
\leanid{TNLean.PEPS.edgeInsertedCoeff}. Together, these declarations represent
the open-edge contraction used when an arbitrary matrix is inserted on the
distinguished edge, as in the argument surrounding Equation
\path{eq:inj_equal_edge} of Ref.~\cite{Molnar2018NormalPEPS}.

The maps \leanid{TNLean.PEPS.edgeMiddleConfigToOpenMiddleConfig} and
\leanid{TNLean.PEPS.edgeOpenMiddleConfigToMiddleConfig} form the finite
equivalence \leanid{TNLean.PEPS.edgeMiddleConfigEquivOpenMiddleConfig}. This
equivalence isolates the reindexing needed for identity insertion. The theorem
\leanid{TNLean.PEPS.edgeMiddleWeight_eq_edgeOpenMiddleWeight} proves the
corresponding middle-tensor reindexing. The diagonal summand is isolated by
\leanid{TNLean.PEPS.edgeInsertedCoeff_identity_diagonal_summand}, while
\leanid{TNLean.PEPS.edgeInsertedCoeff_identity_offDiagonal_summand} proves
pointwise vanishing of every off-diagonal summand. The full specialization is
\leanid{TNLean.PEPS.edgeInsertedCoeff_identity}; its proof reindexes the finite
inserted-boundary sum along the diagonal using
\leanid{TNLean.PEPS.edgeBoundaryConfigEquivDiagonalInsertedBoundaryConfig}.

\section{Formalized Section 3 route}

The following items record the formal realization of the proof in Section~3 of
Ref.~\cite{Molnar2018NormalPEPS}.\footnote{This list and the ledger below are
the status record referenced by the Section~3 docstrings and blueprint
comments.}

\begin{enumerate}
\item
Vertex injectivity passes to the edge-blocked three-site object. The middle
physical index is
\leanid{TNLean.PEPS.EdgeMiddlePhysicalConfig}, and the corresponding middle
contractions are \leanid{TNLean.PEPS.edgeMiddleWeightOn} and
\leanid{TNLean.PEPS.edgeOpenMiddleWeightOn}. The blueprint names the resulting
injectivity statement
\leanid{TNLean.PEPS.EdgeBlockedThreeSiteInjective}.

For a one-vertex region $\{v\}$, the tensor family
\leanid{TNLean.PEPS.singletonRegionTensorFamily} satisfies
\begin{align}
    T_{A,\{v\}}(\eta,\sigma)
        &= A_v(\eta,\sigma).
    \label{eq:peps_injective_ft_section3_route_singleton_region_tensor}
\end{align}
Hence vertex injectivity gives injectivity of the singleton-region tensor.
This base case is formalized by
\leanid{TNLean.PEPS.IsVertexInjective.singletonRegionTensorInjective}. The
passage to the abstract region-injectivity predicate is recorded by
\leanid{TNLean.PEPS.SingletonRegionInjectivityComparison} and
\leanid{TNLean.PEPS.SingletonRegionInjectivityFromVertexInjectivity.ofSingletonComparison}.

The abstract union-closure lemma
\leanid{TNLean.PEPS.RegionInjectivityUnionClosure.finset_injective_of_singletons}
then proves injectivity for every nonempty finite region $R$. In its notation,
\begin{align}
    R
        &= \bigcup_{v\in R}\{v\},
    \label{eq:peps_injective_ft_section3_route_region_singleton_union}\\
    (\forall v\in R,\ \kappa(\{v\}))
        &\Longrightarrow \kappa(R).
    \label{eq:peps_injective_ft_section3_route_region_injectivity_union}
\end{align}

Let the middle region associated with $e=(u,v)$ be
\begin{align}
    M_e
        &= V\setminus\{u,v\}.
    \label{eq:peps_injective_ft_section3_route_middle_region}
\end{align}
The theorem
\path{thm:peps_edgeMiddleTensorInjective_of_singletons} states the corresponding
nonempty-middle consequence. If $M_e\ne\varnothing$ and
\begin{align}
    M_e
        &= \bigcup_{w\in M_e}\{w\},
    \label{eq:peps_injective_ft_section3_route_middle_singleton_union}\\
    (\forall w\in M_e,\ \kappa(\{w\}))
        &\Longrightarrow \kappa(M_e),
    \label{eq:peps_injective_ft_section3_route_middle_region_injectivity}
\end{align}
then the edge-middle comparison gives injectivity of the middle tensor. The
endpoint reduction yields the one-edge theorem
\path{thm:peps_edgeBlockedThreeSiteInjective_of_singletons}; applying it to
every edge with nonempty $M_e$ gives
\path{thm:peps_edgeBlockedThreeSiteInjective_all_of_singletons}.

The graph-cardinality variant
\leanid{TNLean.PEPS.EdgeMiddleRegionInjectivityComparison.%
edgeBlockedThreeSiteInjective_all_two_lt_card}
covers $2<|V|$. For every edge, it uses
\begin{align}
    |M_e|
        &= |V|-2
        > 0,
    \label{eq:peps_injective_ft_section3_route_middle_region_cardinality}
\end{align}
so the nonempty-middle hypothesis is automatic.

Vertex injectivity supplies endpoint injectivity through
\leanid{TNLean.PEPS.IsVertexInjective.edgeBlockedEndpointTensorMaps_injective}
and, uniformly over all edges, through
\leanid{TNLean.PEPS.IsVertexInjective.edgeBlockedEndpointTensorMaps_injective_all}.
The theorem
\leanid{TNLean.PEPS.IsVertexInjective.edgeBlockedThreeSiteInjective_of_middle}
reduces the three-site result to middle-block injectivity, and
\leanid{TNLean.PEPS.IsVertexInjective.edgeBlockedThreeSiteInjective_all_of_middle}
is its all-edge form.

The structure
\leanid{TNLean.PEPS.EdgeMiddleRegionInjectivityComparison} records the
conditional passage from injectivity of the finite region $V\setminus\{u,v\}$
to injectivity of the edge-middle tensor family. Its theorem
\leanid{TNLean.PEPS.EdgeMiddleRegionInjectivityComparison.edgeBlockedThreeSiteInjective}
combines that comparison with vertex injectivity at one edge. The uniform
middle-tensor theorem is
\leanid{TNLean.PEPS.EdgeMiddleRegionInjectivityComparison.edgeMiddleTensorInjective_all},
and the all-edge conditional result is
\leanid{TNLean.PEPS.EdgeMiddleRegionInjectivityComparison.edgeBlockedThreeSiteInjective_all}.

A separate finite-induction route is encoded by
\leanid{TNLean.PEPS.FiniteRegionKernelDescent}. For kernel conditions $K_c(S)$
attached to a coefficient family $c=(c_\rho)_\rho$, the assumptions
\begin{align}
    K_c(S)
        &\Longrightarrow K_c(S\setminus\{j\})
    \label{eq:peps_injective_ft_section3_route_kernel_vertex_deletion}
\end{align}
imply
\begin{align}
    K_c(M_e)
        &\Longrightarrow K_c(\varnothing).
    \label{eq:peps_injective_ft_section3_route_kernel_descent}
\end{align}
The edge-middle instance is
\leanid{TNLean.PEPS.EdgeMiddleKernelDescentData}. A zero relation
\begin{align}
    \sum_\rho c_\rho T_{A,M_e}^{\rho}(\tau)
        &= 0
        \quad\text{for every middle physical index }\tau
    \label{eq:peps_injective_ft_section3_route_middle_zero_relation}
\end{align}
gives $K_c(M_e)$, while $K_c(\varnothing)$ gives
\begin{align}
    c_\rho
        &= 0
        \quad\text{for every }\rho.
    \label{eq:peps_injective_ft_section3_route_terminal_coefficients}
\end{align}
Consequently,
\leanid{TNLean.PEPS.EdgeMiddleKernelDescentData.edgeMiddleTensorInjective}
proves injectivity of $T_{A,M_e}$ from such data, and
\leanid{TNLean.PEPS.IsVertexInjective.edgeBlockedThreeSiteInjective_of_kernelDescent}
combines this result with endpoint injectivity.

Vertex injectivity alone does not imply middle-block injectivity when a virtual
bond has dimension zero. If a complement bond makes
\leanid{TNLean.PEPS.EdgeComplementConfig} empty, then
\leanid{TNLean.PEPS.edgeMiddleTensorFamily_eq_zero_of_isEmpty} proves that the
middle tensor family is identically zero. If the boundary-label index is
nonempty,
\leanid{TNLean.PEPS.not_edgeMiddleTensorInjective_of_isEmpty} proves that this
family is not linearly independent. A four-cycle with a zero-dimensional
middle-to-middle edge gives a mathematical example in which vertex injectivity
still holds vacuously at vertices incident to that edge. Injective PEPS in
Ref.~\cite{Molnar2018NormalPEPS} have nonzero-dimensional virtual bond spaces,
so the formal theorem
\leanid{TNLean.PEPS.IsVertexInjective.edgeBlockedThreeSiteInjective} includes
the hypothesis
\begin{align}
    \forall f,\qquad 0<\dim(f).
    \label{eq:peps_injective_ft_section3_route_positive_bond_dimensions}
\end{align}
The preceding two obstruction lemmas formally verify why this hypothesis is
needed.

The resulting mathematical theorem is the contraction result used in
Ref.~\cite{Molnar2018NormalPEPS}: contracting vertex-injective tensors over a
finite middle region gives an injective blocked tensor. The source proves this
directly from one-sided inverses. The inverse of a contraction of injective
tensors is obtained by contracting their inverses and including the connecting
bond-dimension factor
(\path{Papers/1804.04964/paper_normal.tex:205--250}). The kernel descent above
is a finite-induction route to the same result, not a separate source argument.
The complete transfer is
\leanid{TNLean.PEPS.IsVertexInjective.edgeBlockedThreeSiteInjective}, with the
positive-bond hypothesis.

The proof uses three ingredients. First,
\leanid{TNLean.PEPS.IsVertexInjective.localCoeff_eq_zero_of_contract_zero}
states the one-sided-inverse step: a coefficient family on the local virtual
configurations at a vertex $v$ vanishes if its contraction is the zero physical
vector. Second, the middle tensor is the vertexwise product
\begin{align}
    T_{A,M_e}(\rho,\tau)
        &= \sum_{\zeta}\prod_{v\in M_e}A_v(\zeta,\tau_v),
    \label{eq:peps_injective_ft_section3_route_middle_tensor_product}\\
    \prod_{v\in M_e}A_v
        &= A_j\prod_{v\in M_e\setminus\{j\}}A_v.
    \label{eq:peps_injective_ft_section3_route_middle_vertex_factorization}
\end{align}
Third, the following formal facts instantiate kernel descent:
\begin{enumerate}
\item
The equivalence \leanid{TNLean.PEPS.edgeComplementConfigSplitAt} splits a
complement-bond configuration into the star bonds at $j$, which form the local
virtual configuration at $j$, and the remaining bonds. Its first-component law
agrees with $\textsf{edgeComplementValue}$. In particular, the shared bond
between two middle vertices is recovered from the star factor, so the peeled
coefficient genuinely depends on the local virtual configuration at $j$.

\item
The partial-contraction kernel condition on a finite open region $S$ satisfies
the deletion implication
(\ref{eq:peps_injective_ft_section3_route_kernel_vertex_deletion}). Its proof
fixes the exposed bonds, varies the open physical leg at $j$, and applies
\leanid{TNLean.PEPS.IsVertexInjective.localCoeff_eq_zero_of_contract_zero}.

\item
The terminal implication $K_c(\varnothing)\Rightarrow c=0$ isolates the
residual endpoint boundary labels. Surjectivity of the boundary labelling gives
this implication, and positive bond dimensions ensure that the required labels
exist.
\end{enumerate}
These facts assemble into
\leanid{TNLean.PEPS.edgeMiddleKernelDescentData} and complete the existing
reductions.

\item
The two internal steps of Lemma \path{inj_isomorph} in
Ref.~\cite{Molnar2018NormalPEPS} are formalized for the edge-blocked three-site
MPS. The endpoint realization $X\mapsto O_1,O_2$ is
\path{thm:peps_virtualInsertionPhysicalRealization}: vertex injectivity converts
a one-leg virtual matrix action at either endpoint into a physical operator on
that endpoint. The coefficient-level statement
\path{thm:peps_edgeInsertedCoeff_endpointPhysicalRealization} says that, after
ordinary edge-boundary data are fixed, the endpoint operators supply the
independent inserted-edge index. At the left endpoint, the formal statement
uses the transpose of the inserted matrix because the ordinary boundary
supplies the right bond index.

The converse recovery $O_1,O_2\mapsto W$ is
\leanid{TNLean.PEPS.physical_to_virtual_insertion}.
The theorem \leanid{TNLean.PEPS.resonate_middle_inverted} contracts out the
middle block. The theorem
\leanid{TNLean.PEPS.resonate_invert_right_endpoint} recovers a common matrix
$M$ from the action of $O_1$, while
\leanid{TNLean.PEPS.resonate_invert_left_endpoint} converts the action of $O_2$
into a bond matrix. The theorem
\leanid{TNLean.PEPS.resonate_endpoint_coeff_reconcile} proves equality of these
two matrices, which is the source step $V=W$. Positive virtual bond dimensions
provide the reference residual configurations used to fix $M$.

The local recovery component is
\leanid{TNLean.PEPS.edgeEndpointLocalVirtualOpOfPhysicalOp_eq_of_projected_realization_eq}.
Once the compressed endpoint actions realize the same bond matrix, their local
virtual pullbacks recover that matrix on the shared bond. The conditional
image-preserving theorem
\leanid{TNLean.PEPS.edgePhysicalToVirtualInsertion_of_projected_realization_eq}
records the endpoint conclusion under explicit projected-realization and
image-preservation hypotheses.

\item
Equality of PEPS states is upgraded, after the three-site reduction, to
equality of edge-inserted coefficients. This step produces the matrix-algebra
isomorphism on the chosen bond. The theorem
\leanid{TNLean.PEPS.exists_edgeInsertedCoeff_eq} states that, for every matrix
$X$ inserted into the first blocked PEPS, there is a matrix $Y$ on the
corresponding bond of the second blocked PEPS such that the inserted
coefficients agree at every physical configuration.

The proof uses
\leanid{TNLean.PEPS.edgeRealizationSum_left_eq_sum_edgeBlockedCoeff} and
\leanid{TNLean.PEPS.edgeRealizationSum_right_eq_sum_edgeBlockedCoeff}. These
theorems rewrite the endpoint-realization sums from
\leanid{TNLean.PEPS.edgeInsertedCoeff_eq_sum_left_physicalRealization} and
\leanid{TNLean.PEPS.edgeInsertedCoeff_eq_sum_right_physicalRealization} as
weighted sums of the edge-blocked coefficient. The rewrites use the endpoint
update invariances
\leanid{TNLean.PEPS.edgeOpenMiddleWeight_update_left} and
\leanid{TNLean.PEPS.edgeOpenMiddleWeight_update_right}. Equality of PEPS states
then transfers the sums between the two tensor families through
\leanid{TNLean.PEPS.edgeRealizationSum_left_sameState} and
\leanid{TNLean.PEPS.edgeRealizationSum_right_sameState}. Finally,
\leanid{TNLean.PEPS.physical_to_virtual_insertion} supplies $Y$ for the second
family.

The correspondence is implemented by the explicit map
\leanid{TNLean.PEPS.edgeTransferMatrix}. It composes algebra maps: first insert
$X$ at the right endpoint of the first family and realize it through
\leanid{TNLean.PEPS.edgeRightInsertionOp}; next transfer the resulting physical
operator to the second family using equality of states; finally recover the
matrix realized there. The identity
\leanid{TNLean.PEPS.edgeRightInsertionOp_realizes_edgeTransferMatrix} proves
that this explicit map agrees with the matrix recovered by
\leanid{TNLean.PEPS.physical_to_virtual_insertion}. Because each stage respects
composition,
\leanid{TNLean.PEPS.edgeTransferMatrix_mul} proves multiplicativity.
Theorems \leanid{TNLean.PEPS.edgeTransferMatrix_add} and
\leanid{TNLean.PEPS.edgeTransferMatrix_smul} prove additivity and homogeneity,
and \leanid{TNLean.PEPS.edgeTransferMatrix_edgeInsertedCoeff} proves the
coefficient identity. This construction supplies the multiplicativity that
cannot be read directly from coefficient equality.

Under positive bond dimensions for both tensor families,
\leanid{TNLean.PEPS.isEdgeBlockedInsertionAlgebraIsomorphism} proves that the
correspondence is a $\mathbb{C}$-algebra isomorphism. These hypotheses expose
the nonzero virtual spaces used in Section~3 of
Ref.~\cite{Molnar2018NormalPEPS}, including the reference residual
configurations and the full matrix algebras on the bonds. Injectivity does not
imply them. With a zero-dimensional bond, injectivity and equality of states
can hold vacuously while physical-to-virtual recovery and the gauge conclusion
fail. The theorem-level counterexample is
\leanid{TNLean.PEPS.FundamentalTheoremCounterexample.fundamentalTheoremPEPSStatement_false};
the intermediate recovery failure is
\leanid{TNLean.PEPS.PhysicalToVirtualCounterexample.physical_to_virtual_insertion_statement_false}; and the edge-middle
obstruction is
\leanid{TNLean.PEPS.not_edgeMiddleTensorInjective_of_isEmpty}. Thus the
positive-bond assumptions formalize the source's nonzero virtual bond spaces
rather than an unproved consequence of injectivity.

The final step proves that the edge-inserted coefficient is injective as a
function of the inserted matrix: equality of all physical coefficients
determines that matrix. This gives the unit law
\begin{align}
    \Phi(1)
        &= 1,
    \label{eq:peps_injective_ft_section3_route_transfer_unit}
\end{align}
because identity insertion and $\Phi(1)$ have the same coefficient by equality
of states. It also gives injectivity of the transfer map. Surjectivity follows
from the symmetric construction with the tensor families exchanged. This
matrix injectivity is the inverse-direction content of the resonance inversion
inside \leanid{TNLean.PEPS.physical_to_virtual_insertion}.

\item
The algebra isomorphism is converted into an invertible edge gauge by the
finite-dimensional full-matrix-algebra argument used in the one-dimensional
injective theorem of Ref.~\cite{Molnar2018NormalPEPS}. The theorem
\leanid{TNLean.PEPS.edgeGaugeFromInsertionAlgebraIsomorphism} first identifies
the matrix sizes and then proves that the algebra isomorphism is conjugation by
an invertible matrix. The edge gauges used in the global argument are collected
by \leanid{TNLean.PEPS.exists_edgeGaugeFamily}.

\item
Absorbing all edge gauges into the second tensor family produces modified
tensors $\widetilde B_v$. The formal analog of Equation
\path{eq:inj_equal_edge} of Ref.~\cite{Molnar2018NormalPEPS} is
\leanid{TNLean.PEPS.post_absorption_edge_insertion_equality}. Its conclusion
constructs
\leanid{TNLean.PEPS.PostAbsorptionEdgeInsertionEquality}: inserting the same
arbitrary matrix on the same edge gives equal PEPS coefficients for the first
tensor family and the absorbed second tensor family.

\item
The generalized version of Lemma \path{lem:inj_equal_tensors_2} in
Ref.~\cite{Molnar2018NormalPEPS} is
\leanid{TNLean.PEPS.two_injective_tensor_insertion_comparison}. Equality of all
virtual-bond insertions for two pairs of injective tensors implies reciprocal
scalar proportionality,
\begin{align}
    A_1
        &= \lambda B_1,
    \label{eq:peps_injective_ft_section3_route_first_tensor_scalar}\\
    A_2
        &= \lambda^{-1}B_2.
    \label{eq:peps_injective_ft_section3_route_second_tensor_scalar}
\end{align}
The formal theorem allows nonempty spectator external boundary spaces. The
source lemma is its specialization to one-point spectator spaces.

The coefficientwise rank-one cancellation theorem is
\leanid{TNLean.PEPS.twoBlockReciprocalScalarProportional_of_pointwise_mul_eq}.
Its hypothesis is
\begin{align}
    A_1(\eta_1,\mu,\sigma_1)
    A_2(\eta_2,\nu,\sigma_2)
        &=
    B_1(\eta_1,\mu,\sigma_1)
    B_2(\eta_2,\nu,\sigma_2)
    \quad\text{for all indices}.
    \label{eq:peps_injective_ft_section3_route_pointwise_product_equality}
\end{align}
Injectivity of $A_1$ and $A_2$ then gives
(\ref{eq:peps_injective_ft_section3_route_first_tensor_scalar}) and
(\ref{eq:peps_injective_ft_section3_route_second_tensor_scalar}).

The one-shared-bond separation theorem is
\leanid{TNLean.PEPS.two_injective_tensor_insertion_comparison_singletonBond}.
Insertion of the matrix unit $E_{\mu,\nu}$ extracts the individual product in
(\ref{eq:peps_injective_ft_section3_route_pointwise_product_equality}), after
which the rank-one cancellation theorem applies. This product-separation route
matches the source only in the single-shared-bond case.

For multiple shared bonds, the proof in Ref.~\cite{Molnar2018NormalPEPS} does
not separate the product. It inverts $A_2$, recovers gauges $Z$, $U$, and $W$
on three pairs of legs, and uses compatibility of their identity forms to show
that each gauge is scalar
(\path{Papers/1804.04964/paper_normal.tex:1157--1204}). The scalar conclusion
is \leanid{TNLean.PEPS.threeLeg_residual_forms_scalar}, and the required
inversion and regrouping are contained in
\leanid{TNLean.PEPS.two_injective_tensor_insertion_comparison_core}.

\item
Both injective blocks in the one-vertex-versus-complement comparison are
formalized. The selected-vertex block is
\leanid{TNLean.PEPS.vertexTwoBlock}, with injectivity theorem
\leanid{TNLean.PEPS.isTwoBlockInjective_vertexTwoBlock}. The complement block
\leanid{TNLean.PEPS.complementTwoBlock} contracts $A$ over
$V\setminus\{v\}$ and opens the contraction along
$\textsf{IncidentEdge}(G,v)$ as the shared boundary. Its injectivity theorem,
\leanid{TNLean.PEPS.isTwoBlockInjective_complementTwoBlock}, follows from
vertex injectivity and positive bond dimensions through
\leanid{TNLean.PEPS.regionBlockedTensorInjective_of_isVertexInjective},
specialized to $V\setminus\{v\}$. The boundary-edge and physical-vertex
identifications in
\leanid{TNLean.PEPS.regionBlockedWeight_vertexComplement} identify this
blocked tensor with the vertex-complement tensor. Therefore,
\leanid{TNLean.PEPS.vertexComplementTensorInjective_of_isVertexInjective}
uses the same contraction-of-injective-tensors theorem, not a second
kernel-descent argument.

\item
The final two-tensor comparison blocks one vertex against its complement. In
Ref.~\cite{Molnar2018NormalPEPS}, this gives
\begin{align}
    A_i
        &= \lambda_i\widetilde B_i
        \quad\text{at every vertex }i,
    \label{eq:peps_injective_ft_section3_route_vertex_scalar_comparison}
\end{align}
after which the scalars are incorporated into the local gauges. The theorem
\leanid{TNLean.PEPS.one_vertex_complement_comparison} represents this
specialization. Its formal statement uses abstract two-block notation with
nonempty external boundary spaces; the source comparison follows by taking the
one-vertex and complement blocks after Equation \path{eq:inj_equal_edge} of
Ref.~\cite{Molnar2018NormalPEPS}.

\item
The boundary-insertion argument removes the separate bond-dimension-equality
hypothesis from
\leanid{TNLean.PEPS.fundamentalTheorem_PEPS_of_bondDim}. For each edge,
\leanid{TNLean.PEPS.edgeTransferAlgEquiv} identifies the two full matrix
algebras of virtual insertions. The dimension theorem
\leanid{TNLean.PEPS.bondDim_eq_of_matrixAlgEquiv} then forces equality of the
corresponding bond dimensions. The theorem
\leanid{TNLean.PEPS.fundamentalTheorem_PEPS} gives the gauge-equivalence
conclusion of Ref.~\cite{Molnar2018NormalPEPS} for connected graphs with
positive bond dimensions. The connectivity correction is documented in
\path{docs/paper-gaps/peps_gaugeConsistency_connectivity_gap.tex}.
\end{enumerate}

\section{Source-to-formalization ledger}

The following ledger records the completed correspondence with the Section~3
proof of Ref.~\cite{Molnar2018NormalPEPS}.

\begin{itemize}
\item
For Equation \path{eq:block_to_mps} of
Ref.~\cite{Molnar2018NormalPEPS}, edge blocking is represented by the proved
statement \path{edgeBlockedCoeff_eq_stateCoeff}. The positive-bond injectivity
transfer is
\leanid{TNLean.PEPS.IsVertexInjective.edgeBlockedThreeSiteInjective}. Its
middle physical index is
\leanid{TNLean.PEPS.EdgeMiddlePhysicalConfig}; its named injectivity predicate
is \leanid{TNLean.PEPS.EdgeBlockedThreeSiteInjective}. Endpoint injectivity is
proved both at one edge and uniformly over all edges. The comparison from
edge-middle region injectivity to edge-middle tensor injectivity is also
available uniformly, and its all-edge conditional form is
\leanid{TNLean.PEPS.EdgeMiddleRegionInjectivityComparison.edgeBlockedThreeSiteInjective_all}.

\item
Identity insertion on the chosen bond is
\path{thm:peps_edgeInsertedCoeff_identity}, proved by finite diagonal
reindexing.

\item
The endpoint-realization direction $X\mapsto O_1,O_2$ in Lemma
\path{inj_isomorph} of Ref.~\cite{Molnar2018NormalPEPS} is represented by
\path{thm:peps_virtualInsertionPhysicalRealization} and
\path{thm:peps_edgeInsertedCoeff_endpointPhysicalRealization}.

\item
The recovery direction $O_1,O_2\mapsto W$ in Lemma
\path{inj_isomorph} of Ref.~\cite{Molnar2018NormalPEPS} is
\leanid{TNLean.PEPS.physical_to_virtual_insertion}. Its middle contraction,
endpoint inversions, and matrix reconciliation are respectively
\leanid{TNLean.PEPS.resonate_middle_inverted},
\leanid{TNLean.PEPS.resonate_invert_right_endpoint},
\leanid{TNLean.PEPS.resonate_invert_left_endpoint}, and
\leanid{TNLean.PEPS.resonate_endpoint_coeff_reconcile}. Positive bond
dimensions provide the reference residual configurations. The common-matrix
endpoint recovery is
\leanid{TNLean.PEPS.edgeEndpointLocalVirtualOpOfPhysicalOp_eq_of_projected_realization_eq},
and
\leanid{TNLean.PEPS.edgePhysicalToVirtualInsertion_of_projected_realization_eq}
records the conditional result under explicit image-preservation hypotheses.

\item
The positive-bond matrix-insertion algebra isomorphism is
\leanid{TNLean.PEPS.isEdgeBlockedInsertionAlgebraIsomorphism}. The existence of
a matching matrix is
\leanid{TNLean.PEPS.exists_edgeInsertedCoeff_eq}; its witness is
\leanid{TNLean.PEPS.edgeTransferMatrix}, and its coefficient identity is
\leanid{TNLean.PEPS.edgeTransferMatrix_edgeInsertedCoeff}. The algebra
equivalence includes well-definedness, injectivity, surjectivity, the unit law,
additivity, homogeneity, and multiplicativity. Its positive-bond hypotheses
express the nonzero virtual spaces assumed in Ref.~\cite{Molnar2018NormalPEPS};
they do not follow from injectivity.

\item
The conversion from an algebra equivalence of full matrix algebras to
conjugation by an invertible edge matrix is
\leanid{TNLean.PEPS.edgeGaugeFromInsertionAlgebraIsomorphism}. The resulting
edgewise gauge family is
\leanid{TNLean.PEPS.exists_edgeGaugeFamily}.

\item
Gauge absorption is defined by
\leanid{TNLean.PEPS.absorbEdgeGauges}, and
\leanid{TNLean.PEPS.post_absorption_edge_insertion_equality} applies the
edge-gauge family to the absorbed tensor family.

\item
Equation \path{eq:inj_equal_edge} of Ref.~\cite{Molnar2018NormalPEPS} is
represented by
\leanid{TNLean.PEPS.post_absorption_edge_insertion_equality}, which constructs
\leanid{TNLean.PEPS.PostAbsorptionEdgeInsertionEquality}.

\item
The scalar step in Lemma \path{lem:inj_equal_tensors_2} of
Ref.~\cite{Molnar2018NormalPEPS} uses
\leanid{TNLean.PEPS.twoBlockReciprocalScalarProportional_of_pointwise_mul_eq}
for rank-one product cancellation and
\leanid{TNLean.PEPS.threeLeg_residual_forms_scalar} for the many-bond scalar
forcing.

\item
The complete comparison in Lemma \path{inj_equal_tensors_2} of
Ref.~\cite{Molnar2018NormalPEPS} is
\leanid{TNLean.PEPS.two_injective_tensor_insertion_comparison}. It permits
nonempty external boundary spaces and gives reciprocal scalar proportionality.
The source lemma is its one-point external-boundary specialization.

\item
The one-vertex-versus-complement application is
\leanid{TNLean.PEPS.one_vertex_complement_comparison}, stated in the abstract
two-block formulation with nonempty external boundary spaces.

\item
The theorem
\leanid{TNLean.PEPS.fundamentalTheorem_PEPS_of_bondDim} proves gauge
equivalence after identifying the bond dimensions. The theorem
\leanid{TNLean.PEPS.fundamentalTheorem_PEPS} removes that separate hypothesis
by the matrix-algebra dimension count, and
\leanid{TNLean.PEPS.gauge_unique_mod_edge_scalars} proves uniqueness modulo
edge scalars.
\end{itemize}

\section{Blueprint diagrams}

The blueprint draws geometric steps where they enter the argument and states
intrinsically algebraic steps through their component equations.

\begin{itemize}
\item
The inline edge-blocking reduction shows the chosen edge and the resulting
three-site chain. Following Equation \path{eq:block_to_mps} of
Ref.~\cite{Molnar2018NormalPEPS}, its endpoint regions are $A'_1$ and $A'_2$,
and its complement region is $A'_3$.
% Former TNPEPSThreeSiteInsertionComparison: migrated inline at
% ch24_peps_ft_edge_kernel_gauges.tex; source paper_normal.tex:255--274,565--580;
% both closed-chain panels have boundary (V,P)=(0,3).

\item
The virtual-insertion comparison between the two closed three-site chains
appears inline in Lemma \path{lem:inj_isomorph} of
Ref.~\cite{Molnar2018NormalPEPS}.
% Former TNPEPSInsertionPhysicalRealization: migrated inline at
% ch24_peps_ft_edge_kernel_gauges.tex; source paper_normal.tex:303--352;
% each panel has boundary (V,P)=(0,3).

\item
The realization of a virtual insertion as a physical operation on either
neighboring site is drawn as a three-panel equality.
% Former TNPEPSPhysicalToVirtualInsertion: migrated inline at
% ch24_peps_ft_edge_kernel_gauges.tex; source paper_normal.tex:355--485;
% each panel has boundary (V,P)=(0,3).

\item
The converse implication, in which two neighboring physical operations
determine one virtual operation on their shared bond, is also drawn inline.
% Former TNPEPSEdgeInsertionEquality: migrated inline at
% ch24_peps_ft_edge_kernel_gauges.tex; source paper_normal.tex:1037--1065;
% each five-site panel has boundary (V,P)=(0,5).

\item
After absorption of the edge gauges into the second tensor family, the
five-site comparison with the distinguished $(2,5)$-edge insertion appears
beside Equation \path{eq:inj_equal_edge} of
Ref.~\cite{Molnar2018NormalPEPS}.
% Former TNPEPSTwoInjectiveTensorInsertionComparison: migrated inline at
% ch24_peps_ft_edge_kernel_gauges.tex; source paper_normal.tex:1067--1094;
% the three-shared-bond subcase has boundary (V,P)=(2,2).

\item
The comparison of two injective tensors with a matrix inserted on a shared
virtual bond is drawn in the three-shared-bond subcase, with the external
virtual indices explicit.
% Former TNPEPSTwoInjectiveGaugeScalarReduction: migrated inline at
% ch24_peps_ft_edge_kernel_gauges.tex; source paper_normal.tex:1157--1193;
% all four panels have boundary (V,P)=(3,1).

\item
The four residual-operator forms involving $Z$, $U$, and $W$ appear in the
proof of Lemma \path{lem:inj_equal_tensors_2} of
Ref.~\cite{Molnar2018NormalPEPS}.
% Former TNPEPSEdgeInsertedCoeff: demoted to the exact defining equation at
% ch24_peps_ft_edge_middle.tex. For fixed e, sigma, and X the coefficient has
% boundary (V,P)=(0,0); the former picture's three open legs had the wrong type.

\item
The inserted-edge coefficient is represented by its defining contraction.
After fixing the edge, physical indices, and inserted matrix, the coefficient
is a scalar. A diagram with three open legs would represent a different object.
% Former TNPEPSEdgeGaugeOrientation: demoted to the endpoint-action equations
% in ch24_peps_ft_foundations.tex and ch24_peps_ft_balanced_edge_scalars.tex;
% the two-vertex tensor would have boundary (deg(u)+deg(w)-2,2).

\item
The edge-gauge orientation convention is stated by the two endpoint actions,
not by a separate contraction.
% Former TNPEPSGaugeVertexAction: demoted to the exact component equations in
% ch24_peps_ft_foundations.tex and ch24_peps_ft_balanced_edge_scalars.tex;
% the general vertex tensor has boundary (deg(v),1).

\item
The action of all incident edge gauges on a vertex is stated through its
valence-independent component equation.
% Former TNPEPSGaugeCancellation: demoted to the exact inverse-contraction
% equation in ch24_peps_ft_foundations.tex; its kernel has boundary (V,P)=(2,0).

\item
Cancellation of the endpoint actions on an internal edge is stated as the
corresponding matrix-inverse identity.
% Former TNPEPSBlockedMiddleLocalGaugeFormula: demoted to the exact component
% equation in ch24_peps_ft_edge_kernel_gauges.tex; the general boundary is
% (V,P)=(deg(v),1), not the five-leg boundary of the former picture.

\item
The local gauge obtained from the blocked-middle comparison is stated by its
component equation at an arbitrary vertex.
% Former TNPEPSEdgeGaugeAbsorption: demoted to the exact component equation in
% ch24_peps_ft_edge_kernel_gauges.tex; paper_normal.tex:1212--1314 is
% vertex-specific and the general boundary is (V,P)=(deg(v),1).

\item
Absorption of the incident edge gauges into $B_v$ is stated by the vertexwise
equation defining $\widetilde B_v$.
% Former TNPEPSLocalGaugeExtraction: demoted to the exact component equation in
% ch24_peps_ft_edge_kernel_gauges.tex; no standalone source picture has equal
% boundary types on the two sides.

\item
The local-gauge extraction step is stated by the component equation expressing
$B_v$ as the action of the incident edge matrices on $A_v$.
% Former TNPEPSGlobalConsistency: demoted to the exact component equation in
% ch24_peps_ft_edge_kernel_gauges.tex; paper_normal.tex:1212--1314 gives the
% vertexwise equations rather than a leg-suppressed scalar picture.

\item
Global consistency is stated by one gauge family satisfying the local
component equation at every vertex.
\end{itemize}

For regions $R\subset S=R\cup\{v\}$, the comparison
\begin{align}
    A_R A_v
        &= A_S
        \propto \widetilde B_S
        = \widetilde B_R\widetilde B_v.
    \label{eq:peps_injective_ft_section3_route_normal_region_comparison}
\end{align}
belongs to the normal-region argument in Section~4 of
Ref.~\cite{Molnar2018NormalPEPS}. It is not part of the
one-vertex-versus-complement theorem in Section~3.

The retained inline diagrams use the blueprint's local dot convention while
distinguishing tensor sites, blocked regions, virtual matrix insertions, and
physical operations. Component equations replace separate diagrams when a
picture would hide boundary data or assert a contraction absent from the
source.

\section{Conclusion}

The connected, positive-bond version of the injective PEPS Fundamental Theorem
from Section~3 of Ref.~\cite{Molnar2018NormalPEPS} is formalized. The existence
theorem \leanid{TNLean.PEPS.fundamentalTheorem\_PEPS} and the uniqueness theorem
\leanid{TNLean.PEPS.gauge\_unique\_mod\_edge\_scalars} are proved and
axiom-clean. The proof includes the three-site injectivity transfer, the
matrix-insertion comparison as a $\mathbb{C}$-algebra isomorphism, the edge
gauge construction, and the one-vertex-versus-complement comparison.

The positive-bond-dimension hypothesis records the nonzero virtual bond spaces
assumed in Ref.~\cite{Molnar2018NormalPEPS}. The theorem
\leanid{TNLean.PEPS.FundamentalTheoremCounterexample.fundamentalTheoremPEPSStatement_false}
refutes the conclusion when this hypothesis is dropped. Connectivity is an
explicit correction to the literal source statement and is documented in
\path{docs/paper-gaps/peps_gaugeConsistency_connectivity_gap.tex}. The
counterexample
\leanid{TNLean.PEPS.GaugeConsistencyConnectivityCounterexample.gaugeConsistencyStatement_false} refutes the
gauge-consistency conclusion without connectivity.

% Use the shared bibliography when it is available. The fallback keeps this
% standalone note compilable when the build copies only the input TeX file.
\bibliographystyle{alpha}
\bibliography{references}

\end{document}
